\subsubsection{オブジェクト指向設計}

\term{オブジェクト指向設計}{Object-Oriented Design} は、オブジェクトを中心に設計する方法である。オブジェクトは、データと操作を一体として持つ概念である。設計では、クラス、属性、メソッド、関連、継承、多態性、責務を考える。

初期のオブジェクト指向設計では、名詞からオブジェクトを見つけ、動詞からメソッドを見つける方法が使われた。現在では、単に名詞を拾うだけでなく、責務を中心に設計する考え方が重要である。クラスが何を知り、何を行い、誰と協調するかを考える。

SOLID原則は、クラス設計でよく参照される考え方である。単一責任、開放閉鎖、リスコフ置換、インタフェース分離、依存関係逆転の頭文字である。これらは、変更しやすく理解しやすいオブジェクト指向設計を目指すための指針である。
